% Options for packages loaded elsewhere
% Options for packages loaded elsewhere
\PassOptionsToPackage{unicode}{hyperref}
\PassOptionsToPackage{hyphens}{url}
\PassOptionsToPackage{dvipsnames,svgnames,x11names}{xcolor}
%
\documentclass[
  english,
  11pt,
]{article}
\usepackage{xcolor}
\usepackage[margin=25mm]{geometry}
\usepackage{amsmath,amssymb}
\setcounter{secnumdepth}{-\maxdimen} % remove section numbering
\usepackage{iftex}
\ifPDFTeX
  \usepackage[T1]{fontenc}
  \usepackage[utf8]{inputenc}
  \usepackage{textcomp} % provide euro and other symbols
\else % if luatex or xetex
  \usepackage{unicode-math} % this also loads fontspec
  \defaultfontfeatures{Scale=MatchLowercase}
  \defaultfontfeatures[\rmfamily]{Ligatures=TeX,Scale=1}
\fi
\usepackage{lmodern}
\ifPDFTeX\else
  % xetex/luatex font selection
\fi
% Use upquote if available, for straight quotes in verbatim environments
\IfFileExists{upquote.sty}{\usepackage{upquote}}{}
\IfFileExists{microtype.sty}{% use microtype if available
  \usepackage[]{microtype}
  \UseMicrotypeSet[protrusion]{basicmath} % disable protrusion for tt fonts
}{}
\makeatletter
\@ifundefined{KOMAClassName}{% if non-KOMA class
  \IfFileExists{parskip.sty}{%
    \usepackage{parskip}
  }{% else
    \setlength{\parindent}{0pt}
    \setlength{\parskip}{6pt plus 2pt minus 1pt}}
}{% if KOMA class
  \KOMAoptions{parskip=half}}
\makeatother
% Make \paragraph and \subparagraph free-standing
\makeatletter
\ifx\paragraph\undefined\else
  \let\oldparagraph\paragraph
  \renewcommand{\paragraph}{
    \@ifstar
      \xxxParagraphStar
      \xxxParagraphNoStar
  }
  \newcommand{\xxxParagraphStar}[1]{\oldparagraph*{#1}\mbox{}}
  \newcommand{\xxxParagraphNoStar}[1]{\oldparagraph{#1}\mbox{}}
\fi
\ifx\subparagraph\undefined\else
  \let\oldsubparagraph\subparagraph
  \renewcommand{\subparagraph}{
    \@ifstar
      \xxxSubParagraphStar
      \xxxSubParagraphNoStar
  }
  \newcommand{\xxxSubParagraphStar}[1]{\oldsubparagraph*{#1}\mbox{}}
  \newcommand{\xxxSubParagraphNoStar}[1]{\oldsubparagraph{#1}\mbox{}}
\fi
\makeatother


\usepackage{longtable,booktabs,array}
\newcounter{none} % for unnumbered tables
\usepackage{calc} % for calculating minipage widths
% Correct order of tables after \paragraph or \subparagraph
\usepackage{etoolbox}
\makeatletter
\patchcmd\longtable{\par}{\if@noskipsec\mbox{}\fi\par}{}{}
\makeatother
% Allow footnotes in longtable head/foot
\IfFileExists{footnotehyper.sty}{\usepackage{footnotehyper}}{\usepackage{footnote}}
\makesavenoteenv{longtable}
\usepackage{graphicx}
\makeatletter
\newsavebox\pandoc@box
\newcommand*\pandocbounded[1]{% scales image to fit in text height/width
  \sbox\pandoc@box{#1}%
  \Gscale@div\@tempa{\textheight}{\dimexpr\ht\pandoc@box+\dp\pandoc@box\relax}%
  \Gscale@div\@tempb{\linewidth}{\wd\pandoc@box}%
  \ifdim\@tempb\p@<\@tempa\p@\let\@tempa\@tempb\fi% select the smaller of both
  \ifdim\@tempa\p@<\p@\scalebox{\@tempa}{\usebox\pandoc@box}%
  \else\usebox{\pandoc@box}%
  \fi%
}
% Set default figure placement to htbp
\def\fps@figure{htbp}
\makeatother



\ifLuaTeX
\usepackage[bidi=basic,shorthands=off]{babel}
\else
\usepackage[bidi=default,shorthands=off]{babel}
\fi
\ifLuaTeX
  \usepackage{selnolig} % disable illegal ligatures
\fi


\setlength{\emergencystretch}{3em} % prevent overfull lines

\providecommand{\tightlist}{%
  \setlength{\itemsep}{0pt}\setlength{\parskip}{0pt}}



 


\makeatletter
\@ifpackageloaded{caption}{}{\usepackage{caption}}
\AtBeginDocument{%
\ifdefined\contentsname
  \renewcommand*\contentsname{Table of contents}
\else
  \newcommand\contentsname{Table of contents}
\fi
\ifdefined\listfigurename
  \renewcommand*\listfigurename{List of Figures}
\else
  \newcommand\listfigurename{List of Figures}
\fi
\ifdefined\listtablename
  \renewcommand*\listtablename{List of Tables}
\else
  \newcommand\listtablename{List of Tables}
\fi
\ifdefined\figurename
  \renewcommand*\figurename{Figure}
\else
  \newcommand\figurename{Figure}
\fi
\ifdefined\tablename
  \renewcommand*\tablename{Table}
\else
  \newcommand\tablename{Table}
\fi
}
\@ifpackageloaded{float}{}{\usepackage{float}}
\floatstyle{ruled}
\@ifundefined{c@chapter}{\newfloat{codelisting}{h}{lop}}{\newfloat{codelisting}{h}{lop}[chapter]}
\floatname{codelisting}{Listing}
\newcommand*\listoflistings{\listof{codelisting}{List of Listings}}
\makeatother
\makeatletter
\makeatother
\makeatletter
\@ifpackageloaded{caption}{}{\usepackage{caption}}
\@ifpackageloaded{subcaption}{}{\usepackage{subcaption}}
\makeatother
\usepackage{bookmark}
\IfFileExists{xurl.sty}{\usepackage{xurl}}{} % add URL line breaks if available
\urlstyle{same}
% fallback for those not using the hyperref driver hyperxmp:
\makeatletter
\@ifundefined{xmpquote}{\newcommand{\xmpquote}[1]{#1}}{}
\makeatother
\hypersetup{
  pdftitle={What computer digits can hide},
  pdflang={en},
  colorlinks=true,
  linkcolor={blue},
  filecolor={Maroon},
  citecolor={Blue},
  urlcolor={Blue},
  pdfcreator={LaTeX via pandoc}}


\title{What computer digits can hide}
\usepackage{etoolbox}
\makeatletter
\providecommand{\subtitle}[1]{% add subtitle to \maketitle
  \apptocmd{\@title}{\par {\large #1 \par}}{}{}
}
\makeatother
\subtitle{An optional guide to rounding and sensitive calculations}
\author{}
\date{2026-09-12}
\begin{document}
\maketitle


\href{../paper-map/article.html}{Return to the paper map} · Optional
numerical side article

A computer can perform a calculation exactly as instructed and still
produce a misleading answer. One reason is that it stores only a finite
number of digits.

That matters in this project because the construction repeatedly cancels
quantities. When the desired answer is a small difference between large,
nearly equal contributions, the digits that look least important can
determine the entire result.

We will use subtraction and two elementary equations to separate two
problems: loss of information during calculation and sensitivity already
present in the problem.

\subsection{Subtract first, or round
first?}\label{subtract-first-or-round-first}

Take the numbers 1.2344 and 1.2343. Their difference is 0.0001.

Now keep only four significant digits before subtracting. Both become
1.234, and their stored difference becomes zero.

{\def\LTcaptype{none} % do not increment counter
\begin{longtable}[]{@{}ll@{}}
\toprule\noalign{}
Operation & Result \\
\midrule\noalign{}
\endhead
\bottomrule\noalign{}
\endlastfoot
Subtract the original numbers & 0.0001 \\
Round each number first & 1.234 and 1.234 \\
Subtract the rounded numbers & 0 \\
\end{longtable}
}

The subtraction is not faulty. The needed information was lost before it
began.

\begin{figure}

\centering{

\pandocbounded{\includegraphics[keepaspectratio,alt={Original values 1.2344 and 1.2343 yield 0.0001; rounded values 1.234 and 1.234 yield zero.}]{figures/reader.png}}

}

\caption{\label{fig-rounding}Two nearby numbers remain distinct when
their extra digits are retained. Rounding both to the same stored value
erases their small difference.}

\end{figure}%

The same issue appears when large correction terms are designed to
cancel. A coefficient can look accurate to many digits while still
lacking the particular digits needed for the final small remainder.

\subsection{A better formula can avoid losing the
difference}\label{a-better-formula-can-avoid-losing-the-difference}

Two formulas can mean exactly the same thing in algebra but behave
differently with rounded numbers.

Suppose one formula first builds two nearly equal quantities and
subtracts them. Another calculates their small difference more directly.
The second may preserve information that the first discards.

The supporting study uses such a reformulation for two nearly identical
moment weights. It also uses higher-precision arithmetic where needed.

These are improvements in how we evaluate the same mathematical problem.
They do not alter the intended correction or create new fluid physics.

\subsection{Some problems are sensitive even with perfect
arithmetic}\label{some-problems-are-sensitive-even-with-perfect-arithmetic}

Consider

\[
x+y=2,\qquad x+1.000001y=2.000001.
\]

Subtract the first equation from the second. We get
\(0.000001y=0.000001\), so \(y=1\) and \(x=1\).

Now change only the second right-hand side to 2.000002. The subtraction
gives \(y=2\), so \(x=0\).

A tiny change in the supplied number caused a large change in the
answer. No rounding error is needed for that result. The two equations
ask nearly the same question, so their small difference carries almost
all the information about how to distinguish \(x\) and \(y\).

This is called \textbf{poor conditioning}. It is a property of the
problem and its chosen variables. More accurate arithmetic can reveal
the sensitivity; it cannot remove genuine uncertainty in the input data.

\subsection{Keep the extra digits all the way
through}\label{keep-the-extra-digits-all-the-way-through}

Suppose a calculation uses many digits to determine accurate
coefficients, then saves them in an ordinary lower-precision format. The
later cancellation may lose the benefit of the careful calculation.

That is like measuring a length to the nearest micrometre and writing it
down only to the nearest centimetre.

For a useful accuracy claim, we must track the full path: input
accuracy, calculation, storage, and the final quantity being checked.

The same lesson applies to coordinates. A finely spaced sampling grid is
ineffective if its points round to identical numbers before evaluation.
The \href{../d6_global_axis/article.html}{narrow-band study} provides an
example from the inner-profile calculations.

\subsection{What our experiment
tested}\label{what-our-experiment-tested}

The numerical study examined a finite moment-repair problem from the
companion material. As a parameter changes, two of its weighted
constraints become increasingly alike.

We compared equivalent formulas, retained different numbers of digits,
and changed the supplied discrepancies slightly. These experiments
separated avoidable numerical loss from the sensitivity of the
underlying solve.

This is a test of an algebraic component. It does not measure whether a
fluid flow is dynamically stable, and it is not a numerical check of the
whole proof.

\subsection{Where this fits in the bigger
picture}\label{where-this-fits-in-the-bigger-picture}

\textbf{Computer accuracy and mathematical validity are different
questions, and both need attention.} A computation can illustrate a
proof block only if it preserves the information that the block uses.

This is an optional lesson about the companion's numerical methods, not
an additional mechanism in the paper. Return to the
\href{../paper-map/article.html}{paper map} for the main reading route.

The \href{technical-article.html}{technical article},
\href{../../code/d7_conditioning/study.ipynb}{notebook}, and
\href{source-notes.md}{source notes} record the actual moment weights,
equivalent formulation, and precision experiments.




\end{document}
